% !TEX root = sbi.tex

%==============================================================================
\chapter{Introduction}
\label{sec:intro}
%==============================================================================
The questions on how did the universe came to be, why is the universe as such, and what is the universe heading towards falls upon cosmology to answer. Cosmology is the study of the universe on large scales in an attempt to answer the latter questions {\color{green} cite maybe Peebles}. This pursuit started as a data-starved research {\color{green} cite} and transformed into a data-rich endevour as a result of a platitutde of imaging surveys such a-s the LSST Vera Rubin observatory, Euclid, and Nancy Roman {\color{green} cite}. The qualitative questions, with which this paragraph started,< transcribe in the statistical sense to be what is the cosmological model that best describes the universe on throughout length scales and energy scales. 

The first attempt to describe the universe as a whole came through the lens of the general theory of relativity where Robertson and Walker  {\color{green} cite} built on the work of Friedmann and Lemaître to describe geometry of the universe through what is now known as the Friedmann-Lemaître-Robertson-Walker (FLRW) metric. We discuss this metric further in \ref{sec:FLRW}. This metric was largely motivated by the observations of Edwin Hubble who saw that galaxy further away from the Earth are moving faster than galaxies closer to Earth {\color{green} cite}. Figure shows the plot from Hubble's work in 1929. 

Further observations of the magnitude of distant galaxies have shown that the universe is expanding {\color{green} cite}. The general theory of relativity attributes this expansion to the energy density content of the universe. Therefore, the question of what is driving the expansion of the universe and what are the componenets that contribute to the energy budget of the universe represent a crucial question in cosmology for which multiple cosmological models have been suggested \citep{di2025cosmoverse}. In fact, the cosmological model has been motivated by both cosmological and astrophysical observation. Some of these observation are the measurement of galaxy rotation curves by Vera Rubin {\color{green} cite} that hinted at an additional matter component that does not interact with light and the discovery of the cosmic microwave background (CMB) radiation that hinted at hot and homogeneous early universe.{\color{green} cite}. Furthermore,early observaiton of clusters of galaxies have shown that the spatial distribution of galaxy deviates from randomness. This was shown by Peebles {\color{green} cite}. By calculating the two point correlation function, there was clearly a signal in the way that galaxies are distributed which points out to an underlying mechanism in the formation of these galaxies \citep{Totsuji_Kihara_1969}. The development of obsrevational techniques and instrumentation allowed the launch of the Center for Astrophysics Redshift Survey \citep{Davis_Huchra_Latham_Tonry_1982} which showed that galaxy positions are correlated on large scales and that there is a structure in the universe \citep{Davis_Peebles_1983}.

In the early 1990s, cosmology experienced a paradigm shift with the launch of the Cosmic Background Explorer (COBE) sattelite \citep{Boggess_Mather_etal._1992}. The COBE mission provided the first detection of primordial CMB anisotropies and further confirmed that the universe is flat \citep{Smoot_1999}. In the context of the FLRW matric, the flatness of the universe requires a total sum of the relative density parameters to be 1 (within error bounds). Therefore, in addition to the matter component (dark and baryons), there should be a third component that can account for the missing energy density. Using supernova data from the Hubble Space Telescope \citep{Bahcall_Spitzer_1982}, \citet{Riess_etal._1998} found a strong evidence for an expanding universe universe with negative acceleration and supported the findings of the COBE mission for the existance of additional density component different from matter which is driving the expansion of the universe. Hereafter, the term 'dark energy was coined to refer to this component \citep{Turner_2001}. Wiht this, the $\Lambda$CDM model was born, and the efforts to obtain a precise estimates of the cosmological parameters intensified \citep{Bennett_etal._2003, Pope_2004}.

While early experiements relied on CMB, galaxy number counts, or supernova to infer cosmological parameters, gravitational lensing was in development as a tool to constrain the cosmological parameters. Gravitational lensing is the distortions of light geodesics by gravity when light travels from a distant source in the universe towards an observer on the Earth \citep{Schneider_Ehlers_Falco}. On the level, of a galaxy, gravitional lensing effects manifest themselves as a distortion of the shape of the galaxy and a change in the angular position with respect to a particular observer. \citet{Kaiser_Squires_1993} showed that these distortions could be used to obtain a map of the integrated mass density along the line of sight. The use of galaxy shape distortions for cosmology was a direct result of the Kairser-Squire inversion formula where \citet{seitz1996steps} realised that the expectation value of galaxy shape is an unbiased estimate of the shear in the weak field limit. In principle, it is possible to relate the power spectrum of galaxy shapes with the matter power spectrum, offering a concrete method of probing cosmology by studying the correlations between galaxy shapes \citep{Bartelmann_Schneider_2001}. This new probe of the large scale structure is called cosmic shear and was used by \citet{Maoli_VanWaerbeke_2001} to obtain constraints on the cosmological parameters. These constraints were not as tight as the ones obtained by COBE, however, it paved the road towards using cosmic shear to constrain the cosmological model. The advantage of cosmic shear as an integrated quantitity along the line of sight mean that it is sensitive to structure growth in the late universe. Therefore, it represent a very good tool to constrain the equation of state for dark energy \citep{Hu_Tegmark_1999, Jain_Taylor_2003}.

After that, multiple galaxy surveys used cosmic shear for cosmilogical parameter inference, and numerous statistical tools were developped to deal with the increasing data size \citep{Schneider_VanWaerbeke_Jain_Kruse_1998, Schneider_Eifler_Krause_2010}. Likewise, with increasing data size, the noise level in a cosmological signal increases and the need to mititgate the impact of the systematic and measreument uncertainties is imperative. In this work, we focus on tackling on of the most crucial systematics, namely baryonic feedback on the large scale structure.

Baryonic feedback are the set of baryonic physics processes that impact the distribution of the matter \citep{Rudd_Zentner_Kravtsov_2008}. These processes include, but not limited to, supernova feedback, activate galactic nuclei feedback, and star formation. By analying hydrodynamical simulations,  \citet{van2011effects} realised that baryonic feedback affect the matter distribution in the universe at different scales, meaning that their impact should be modelled and integrated . However, the theoretical formalism of structure growth, and therefore the matter power spectrum assumes that structure growth is only driven by gravitional effects, ignoring baryonic feedback \citep{Dodelson_Schmidt_2026}. These uncertainties naturally tunel down to any cosmic shear statistics since it is related to the matter power spectrum. 
To mitigate the impact of baryonic, 

To estimate the angular power psectrum up to small angular scales ($\ell$ \~ 5000), the matter power spectrum should be accuretly modeled up to an equivilant small scales ($k ~ 10 ,/ hMp>^{-1}$)

One of trhe   





To unleash the full potential of cosmic shear 
The theoretical formalism of cosmic shear assumes that structure growth is only driven by gravitional effects. However, hydrodynamical simulations has shown that baryonic feedback affect the matter distribution in the universe at different scales \citet{van2011effects}, and ignoring the impact of baryonic physics leads to biased estimate of the cosmological parameters in weak lensing analysis. The problem of baryonic phyiscs is very unique to using cosmic shear as a probe and multiple methods have been suggested to mitigate their impact.

For example, statistical scale cuts \citep{krause2021dark}, theoretical error covariance (m), SZ effect, 

Speak about the FLRW metric, Hubble flow, energy content, impact of the Hubble flow and initial conditions on 
structure formation and evolution. Dark matter and traces of dark matter. Weak lensing, baryonic feedback
FRBs, simulation based inference.

\begin{itemize}
  \item Large Scale Structure
  \item Lensing
  \item Weak gravitational Lensing
  \item Precision Cosmology \& tensions
  \item Baryonic feedback
  \item FRBs
  \item systematics of FRBs
\end{itemize}


% Print the bibliography starting on the same page at the end of the chapter.
% \printbibliography[heading=subbibliography]
